\documentclass[11pt,a4paper]{article}

%========================= FONTES E CODIFICAÇÃO =========================
\usepackage[T1]{fontenc}
\usepackage[utf8]{inputenc}
\usepackage{mathptmx}
\usepackage[scaled=0.88]{helvet}
\usepackage{courier}
\usepackage{microtype}

\hyphenpenalty=9000
\exhyphenpenalty=9000
\tolerance=2500
\emergencystretch=3em

%========================= GEOMETRIA =========================
\usepackage[a4paper,top=2.6cm,bottom=2.6cm,left=2.7cm,right=2.7cm,headheight=14pt]{geometry}

%========================= PACOTES =========================
\usepackage{xcolor}
\usepackage{booktabs}
\usepackage{tabularx}
\usepackage{array}
\usepackage{enumitem}
\usepackage{titlesec}
\usepackage{fancyhdr}
\usepackage{siunitx}
\usepackage[most]{tcolorbox}
\usepackage[hidelinks,pdfusetitle]{hyperref}
% Resolve \ref{} to plano-port-max10.tex's section numbers automatically
% (reads plano-port-max10.aux -- compile that document at least once first,
% and again after any section renumbering there, before recompiling this one).
\usepackage{xr}
\externaldocument{plano-port-max10}

%========================= CORES (mesma paleta do plano principal) =========================
\definecolor{lescblue}{HTML}{1F3A5F}
\definecolor{lescaccent}{HTML}{2E6DA4}
\definecolor{lescgray}{HTML}{4A4A4A}
\definecolor{lesclight}{HTML}{EEF3F8}
\definecolor{okcolor}{HTML}{2F5D3A}
\definecolor{oklight}{HTML}{EDF4EE}
\definecolor{rulegray}{HTML}{B8C4D0}

\titleformat{\section}
  {\normalfont\Large\bfseries\color{lescblue}}
  {\thesection}{0.7em}{}[{\vspace{2pt}\color{rulegray}\titlerule[0.8pt]}]
\titleformat{\subsection}
  {\normalfont\large\bfseries\color{lescaccent}}
  {\thesubsection}{0.6em}{}
\titlespacing*{\section}{0pt}{16pt}{8pt}
\titlespacing*{\subsection}{0pt}{12pt}{5pt}

\pagestyle{fancy}
\fancyhf{}
\renewcommand{\headrulewidth}{0.4pt}
\renewcommand{\footrulewidth}{0pt}
\fancyhead[L]{\footnotesize\color{lescgray}Nota de direcionamento --- limite de detectabilidade}
\fancyhead[R]{\footnotesize\color{lescgray}LESC/UFC}
\fancyfoot[C]{\footnotesize\color{lescgray}\thepage}

\setlist[itemize]{leftmargin=1.3em,itemsep=2pt,topsep=4pt,parsep=0pt}
\setlist[enumerate]{leftmargin=1.6em,itemsep=3pt,topsep=4pt,parsep=0pt}

\newtcolorbox{keybox}[1][]{
  enhanced, breakable,
  colback=lesclight, colframe=lescaccent,
  boxrule=0.4pt, left=8pt, right=8pt, top=7pt, bottom=7pt,
  arc=1.5pt,
  fonttitle=\bfseries\small\color{white},
  coltitle=white, colbacktitle=lescaccent,
  attach boxed title to top left={xshift=6pt,yshift=-3pt},
  boxed title style={arc=1pt,boxrule=0pt},
  #1}

\newtcolorbox{okbox}[1][]{
  enhanced, breakable,
  colback=oklight, colframe=okcolor,
  boxrule=0.4pt, left=8pt, right=8pt, top=7pt, bottom=7pt,
  arc=1.5pt,
  fonttitle=\bfseries\small\color{white},
  coltitle=white, colbacktitle=okcolor,
  attach boxed title to top left={xshift=6pt,yshift=-3pt},
  boxed title style={arc=1pt,boxrule=0pt},
  #1}

\renewcommand{\arraystretch}{1.2}
\newcolumntype{L}[1]{>{\raggedright\arraybackslash}p{#1}}

\title{Nota de Direcionamento --- O Estudo de Limite de Detectabilidade como Objetivo Primário}
\author{LESC --- Laboratório de Engenharia de Sistemas de Computação, UFC}
\date{}

\begin{document}
\maketitle
\thispagestyle{fancy}

\noindent\textbf{Relação com o plano principal.} Este documento é um complemento a
\texttt{plano-port-max10.tex/pdf}, no mesmo espírito dos artefatos de 2 a 4 páginas
que a Fase~0 exige de cada trilha de estudo (Seção~\ref{sec:fase0-metodo} do plano principal). Não
substitui nenhuma seção do plano; propõe uma mudança de \emph{enquadramento} do
objetivo do projeto, com efeito concreto e limitado sobre um pequeno número de
seções --- listadas ao final.

\section{A proposta em uma frase}

Declarar, desde o dia um, que o objetivo primário do projeto é \textbf{caracterizar
com rigor o limite de detectabilidade de um sensor de \emph{slack} portado entre
fabricante e nó tecnológico}, e tratar a recuperação de resolução por Vernier de
dois PLLs --- hoje a tarefa de maior valor da frente F1 --- como um resultado
\emph{adicional} que, se bem-sucedido, eleva a contribuição, mas cuja falha não
esvazia o semestre.

\section{Por que propor isto agora}

O próprio plano principal já contém o argumento --- só não o usa como ponto de
partida. Três passagens, lidas em conjunto, sustentam a proposta:

\begin{itemize}
  \item A Seção~\ref{sec:propagacao-perda} do plano principal deriva, por cálculo, uma perda de resolução
        de $4{,}86\times$ entre o passo de fase do MMCM e o do PLL do MAX~10 ---
        \emph{antes} de qualquer tentativa de mitigação. Esse número já é, por si
        só, um resultado científico: ninguém publicou a penalidade de portabilidade
        de um sensor de \emph{slack} entre fabricantes com dados de \emph{hardware}
        real.
  \item A Seção~\ref{sec:gateA} (Gate~A) já prevê o desfecho ``Vernier inviável, passo nativo
        funcionando'' como aceitável, reposicionando o projeto como ``estudo de
        limite de detectabilidade de sensores de \emph{slack} em nó tecnológico
        maior'' --- \emph{e afirma explicitamente que isso continua sendo resultado
        publicável}.
  \item A Seção~\ref{sec:resultado-desfavoravel} (``Resultado desfavorável bem medido é resultado'') já instrui
        o orientador a comunicar isso à equipe no dia um, precisamente para evitar
        que, sob pressão de prazo, a equipe tente forçar um resultado favorável.
\end{itemize}

Ou seja: o plano principal já trata o estudo de limite de detectabilidade como
\emph{desfecho aceitável de contingência}. A proposta aqui é promovê-lo a
\emph{objetivo declarado desde o início}, não a resultado de consolação decidido na
semana~8.

\subsection{O que isso muda na prática}
\label{sec:pratica}

\begin{keybox}[title={A diferença não é técnica, é de sequenciamento}]
Em nenhum dos dois enquadramentos o trabalho de engenharia muda --- a frente F1
ainda tenta o Vernier, o Gate~A ainda acontece na semana~8. O que muda é o que
\emph{depende} do sucesso do Vernier. Hoje, formalmente, o sucesso do projeto
depende do Gate~A ter um desfecho favorável. Com a proposta, o sucesso do projeto
depende apenas da caracterização estar bem feita --- o Gate~A decide qual
\emph{versão} do resultado é publicada, não se há resultado.
\end{keybox}

\begin{itemize}
  \item \textbf{Risco do semestre.} Hoje, o pior desfecho documentado no registro de
        riscos (Seção~\ref{sec:riscos}, plano principal) é ``sinal de envelhecimento permanece
        abaixo do piso de ruído'' --- um risco que existe \emph{independente} do
        Vernier. Ancorar o sucesso do projeto na caracterização, e não na
        recuperação de resolução, remove uma dependência de risco (falha do
        Vernier) da definição de sucesso, sem removê-la do escopo de trabalho.
  \item \textbf{Redação da produção científica.} A tarefa~6 da frente F6 (``Produção
        escrita'') passa a ser redigida, desde a primeira versão, em uma forma que
        sobrevive a qualquer desfecho do Gate~A: \emph{``Caracterizamos o limite de
        detectabilidade alcançável para um sensor de \emph{slack} portado entre
        fronteiras de fabricante e de nó tecnológico, e adicionalmente reportamos
        uma tentativa de recuperação de resolução por batimento Vernier de dois
        PLLs, com eficácia caracterizada e limitada por \emph{jitter}.''} Se o
        Vernier funcionar, a segunda oração cresce; se não funcionar, ela encolhe
        para um parágrafo de tentativa caracterizada --- nunca precisa ser reescrita
        do zero.
  \item \textbf{Critério de aceite da frente F6.} Não precisa mais aguardar o
        desfecho do Gate~A para ser considerado satisfeito. Assim que
        $\sigma_\text{resid}$ do passo nativo do PLL estiver medido (o que acontece
        \emph{antes} do Gate~A, como pré-requisito dele --- Seção~\ref{sec:gateA}, evidência~1),
        o critério de aceite de F6 já pode ser declarado cumprido para a versão
        ``limite de detectabilidade''. O Vernier, se bem-sucedido, adiciona uma
        segunda medida ao mesmo critério, não o substitui.
  \item \textbf{Comunicação com a equipe.} Em vez de ``o Gate~A decide se o projeto
        continua sendo uma migração ou vira outra coisa'' (redação atual da
        Seção~\ref{sec:gateA}), a mensagem passa a ser ``o projeto já é, desde o início, uma
        caracterização de limite de detectabilidade; o Gate~A decide se ganhamos
        também uma técnica de mitigação validada''. A diferença psicológica para uma
        equipe em formação, sob pressão de prazo no fim do semestre, não é
        pequena --- é exatamente o tipo de situação que a
        Seção~\ref{sec:resultado-desfavoravel} do plano principal já identifica
        como modo de falha recorrente.
\end{itemize}

\section{O que \emph{não} muda}

Esta proposta não altera nenhuma tarefa de engenharia, nenhum cronograma, nenhuma
alocação de integrante. F1 continua investigando o Vernier com a mesma prioridade;
o Gate~A continua na semana~8 com os mesmos critérios de evidência. A mudança é
inteiramente de enquadramento e de ordem de redação --- baixo custo, baixo risco de
implementação, e reversível a qualquer momento caso a equipe prefira o enquadramento
original.

\section{Extensão natural, se adotada: comparação de três nós}

Com este enquadramento, uma extensão de custo marginal quase nulo se torna
disponível para a frente F6 (Seção~\ref{sec:F6} do plano principal, tarefa~4 ``Comparação
sistemática entre plataformas''): o laboratório já possui, ou terá em breve, dados
de sensor de \emph{slack} com metodologia consistente em \textbf{três} nós/fabricantes
--- UltraScale+ customizado (\SI{16}{\nano\meter} FinFET), Artix-7 (\SI{28}{\nano\meter}
planar) e MAX~10 (\SI{55}{\nano\meter}). Um gráfico de
$\sigma_\text{resid}$ (ou do degrau detectável $3\sigma$) em função do nó, com os
três pontos, é uma figura adicional de baixo esforço que fortalece exatamente o
resultado que esta nota propõe como objetivo primário --- vale considerar incluí-la
no rascunho de artigo da tarefa~6, mesmo que como figura secundária.

\section{Seções do plano principal afetadas, se a proposta for adotada}

\begin{center}
\begin{tabularx}{\textwidth}{@{}L{3.6cm} X@{}}
\toprule
\textbf{Seção (plano principal)} & \textbf{Ajuste sugerido} \\
\midrule
Seção~\ref{sec:objetivo} Objetivo &
Adicionar uma frase declarando a caracterização de limite de detectabilidade como
objetivo primário e a recuperação por Vernier como objetivo secundário. \\
Seção~\ref{sec:gateA} Gate A, ``Desfechos possíveis'' &
Reescrever a lista de três desfechos como duas variantes do mesmo resultado
(caracterização sozinha vs.\ caracterização + mitigação validada), não como
bifurcação de sucesso/fracasso do projeto. \\
Seção~\ref{sec:F6} F6, tarefa 6 (Produção escrita) &
Redigir o rascunho de artigo desde a primeira versão na forma dupla descrita na
Seção~\ref{sec:pratica} desta nota. \\
Seção~\ref{sec:produtos} Produtos esperados &
Nenhuma mudança de conteúdo; ajustar apenas a descrição do produto ``Submissão a
conferência'' para refletir o enquadramento primário. \\
\bottomrule
\end{tabularx}
\end{center}

\end{document}
